\documentclass[ecta,nameyear,draft]{econsocart}
\usepackage{amsmath,amsfonts,amssymb,amsthm,booktabs,array,longtable}
\usepackage[T1]{fontenc}
\usepackage{mathpazo,microtype,needspace,xurl}
\usepackage{xr-hyper}
\RequirePackage[colorlinks,citecolor=blue,linkcolor=blue,urlcolor=blue]{hyperref}
\startlocaldefs
\newtheorem{theorem}{Theorem}
\newtheorem{proposition}[theorem]{Proposition}
\newtheorem{lemma}[theorem]{Lemma}
\newtheorem{corollary}[theorem]{Corollary}
\newcommand{\E}{\mathbb E}
\newcommand{\R}{\mathbb R}
\endlocaldefs
\setlength{\emergencystretch}{3em}
\makeatletter\def\copyright@text{Prepared for referee review}\makeatother

\setcounter{secnumdepth}{2}

% Stable current-document identity for clean typesetting reproduction.
\ifdefined\pdfinfoomitdate\pdfinfoomitdate=1\fi
\ifdefined\pdfsuppressptexinfo\pdfsuppressptexinfo=15\fi
\ifdefined\pdftrailerid\pdftrailerid{}\fi

\externaldocument[][nocite]{build/supp}[supp.pdf]
\externaldocument[][nocite]{build/appendix}[appendix.pdf]
\begin{document}
\begin{frontmatter}
\title{Neural Bellman Operators}
\runtitle{Neural Bellman Operators}
\begin{aug}
\author[id=au1,addressref={add1}]{\fnms{Qian}~\snm{QI}\ead[label=e92ECTA]{qiqian@pku.edu.cn}}
\address[id=add1]{Peking University}
\end{aug}
\begin{abstract}
This paper develops Neural Bellman Operators for continuation learning and feasible decisions under the original recursive economic objective. Centered tail-switch bounds connect nonlinear-risk continuation error to constrained action loss. An active-region Bellman compiler and outward witnesses separate numerical proposals from certificate authority. Completed-incumbent safeguards establish actual policy improvement; a remaining-gain rule stops further computation with a verified bound on foregone improvement. We prove its pointwise and recursive accounts, alongside retained own-future construction and information results. A separately delimited BSDE verification theorem connects the framework to recursive utility. Two complete original-investment comparisons retain 10,240 primary decisions and all fresh-process timing records. At the prespecified gain budget, 2,386 of 3,072 extra trial pairs are omitted, with actual original-policy cost sacrifices reported. Strong conventional controls receive identical information and algorithms; equivalent ReLU and hinge geometry gives identical policies. Historical adverse evidence remains visible. Neural-exclusive competitiveness and calibrated adoption require their own further evidence.
\end{abstract}
\begin{keyword}
\kwd{Dynamic programming}\kwd{Neural approximation}\kwd{Recursive risk}
\kwd{Policy qualification}\kwd{Computational economics}
\end{keyword}
\end{frontmatter}
\section{Introduction}\label{sec:intro110}
Dynamic investment under risk requires a continuation value, a feasible action and a decision about how much computation to purchase. These objects cannot be assessed independently. A more accurately fitted continuation can move a nearly binding action in the wrong direction; a better proposed action can still require an expensive original-model certificate; and a certified policy improvement can be too small to justify another round of computation. Neural Bellman Operators (NBOs) study this joint problem by constructing continuation and policy maps, evaluating their own future decisions, and qualifying their loss under the original economic law.

The economic question is when an acquired neural continuation helps deliver original-objective decisions at a competitive complete cost. The comparison is with conventional representations given the same information, feasible actions and computational components. A theorem whose validity is representation neutral is not evidence that neural networks alone can implement it. Equally, a negative result for one fitting loss does not exhaust what an action-aligned continuation and policy construction can achieve.

We develop an original-law chain from continuation accuracy to feasible decisions and paid stopping. The retained own-future NBO construction returns an exactly realizable neural continuation and a measurable feasible actor without assuming successful nonconvex training. Reference-centered recursive-risk bounds distinguish action-dependent errors from irrelevant level shifts. An original-law supporting quadratic then bounds the gap of the action actually proposed, including constrained stationarity and numerical witness slack. All these accounts refer to the original Bellman operator, not to an easier fitted economy.

Nonlinear risk creates an additional obstruction. Two continuation outcomes can have identical action derivatives at every innovation and still have different risk-optimal actions because their maximizing tails differ. A centered quantile-band theorem bounds this discrepancy using value error, directional error and the oscillation of the action gradient near the tail boundary. Its exact continuous-law diagnostic has zero directional approximation error but original action loss $1/4$. The conditions under which the bound is differentiable, the treatment of boundaries and the separate allowance for an imperfect implemented future policy are explicit. A finite training sample is not a certificate of its uniform margin assumptions.

The second contribution concerns the economic return to further computation. We first derive an active-region terminal Bellman oracle from the original investment primitives. Its action derivative is a strictly increasing cubic on at most four regions. Formula-derived actions are independently qualified by the original outward witness; numerical inversion never acquires certificate authority. This provides a strong semi-analytic control to every representation. We then retain a completed qualified incumbent and accept a new action only when its policy-cost upper bound lies strictly below the incumbent's lower bound. A recursive gain theorem states when such local inequalities support a uniform policy-improvement conclusion.

The current algorithmic addition uses the same original-optimum lower bound to cap \emph{remaining} policy improvement. When that cap is below a predeclared gain budget, a state stops paying for further support and trial evaluations. Qualification targets remain unchanged. We prove the pointwise foregone-gain guarantee, its finite-catalogue account, and its recursive composition under the required comparator-future residual. The stopping rule does not pretend to be globally optimal metareasoning, nor does it assign an invented price to economic welfare. It makes an implementable, testable trade-off between actual original-policy cost and complete computational work.

The numerical comparison preserves the original two-date continuous-risk investment economy. Eight proposal arms share 256 paid terminal-value labels: an online original-model service, rank-six quadratic, ReLU and coordinate-spline continuations, an equivalent conventional hinge, full quadratic and cubic polynomials, and bilinear tensor interpolation. An unchanged-source reproduction retains all 4,096 decisions in the original four-service comparison. The new experiment retains all 6,144 decisions in six services, including the original controls, the primary remaining-gain budget $2^{-16}$ and all three prespecified sensitivities. Both campaigns keep seven sequential fresh-process timing repetitions, all candidate and action rejections, and separate acquisition, audit, reference and service charges. These are deterministic catalogues, not new population samples.

The evidence improves the computation--decision connection without manufacturing a neural victory. At the primary budget, the new rule omits 2,386 of 3,072 possible extra support/trial pairs across the declared arms and risks. For ReLU at risk level $1/2$, original terminal-risk queries fall from 94,737 to 35,832, and the median complete 32-state service clock falls from about $0.4074$ to $0.2074$ seconds. Its mean original cost rises relative to the full safeguard by a certified interval contained in $[1.95,3.34]\times10^{-7}$; that foregone improvement is reported, not hidden as equal accuracy. The budgeted service still strictly improves its completed plain incumbent. Conventional methods receive the same stopping benefit, and ReLU and equivalent hinges return exactly identical policy and trial records.

\paragraph{Relation to established methods and evidence.}
Risk-averse dynamic programming and policy iteration already provide the recursive foundation \citep{ruszczynski2010risk,ruszczynski2014erratum}. Inexact-oracle optimization separates computation from its accuracy account \citep{DevolderGlineurNesterov2014}; decision-focused learning and CVaR differentiation provide related acquisition and gradient principles \citep{ElmachtoubGrigas2022,Tamar2015}. Our results specify their connection to original-law supporting witnesses, tail switching, completed feasible policies, risk-distorted gains and charged resolution. The retained information bounds have explicit oracle restrictions; a formula-derived compiler is not a counterexample to a point-value lower bound. Recent static trajectory-CVaR work \citep{LuoDelage2026} optimizes a different recursive object from the nested AVaR preference here. Its numerical improvements cannot be transferred between those objectives.

The original NBO identity includes rigorous connections to recursive utility and stochastic control. A self-contained Brownian Markov BSDE verification result in the supplement states one admissible bridge, including well-posedness, the critic computational graph and terminal or stopping-boundary errors. It does not assume equivalence with a general coupled NDU/FBSDE or claim that the investment catalogue evaluates that different utility model.

The historical quadratic complete-work winner, the increase from 644 to 757 first-proposal rejections despite improvement of all 36 residual spans, and zero certified learned adoption over quadratic are retained. Current strict policy improvements and work savings do not erase those observations. Nor do they establish a neural-exclusive competitive region, calibrated adoption or an unexecuted longer-horizon tournament. The main article concentrates on the original economic question, current mathematical results and complete comparisons; all earlier theorems, proofs, negative campaigns and source identities remain in the Technical Supplement, Development Appendix and source archive.

\section{The economic operator and its policy account}\label{sec:core49}
Consider dates $t=0,\ldots,T$ and an invariant compact state domain $X\subset\R^d$. At date $t$, the action belongs to a nonempty compact set $A_t(x)$. The transition is $F_t(x,a,Z_{t+1})$, the current cost is $c_t(x,a)$, and the terminal cost is $g(x)$. Write $B_0=1$ and $B_t=\prod_{j<t}\beta_j$, with $\beta_j>0$. An implemented policy includes its observation rule, action selection and feasible numerical realization, not just its stored parameters.

For a continuation $v$, define
\begin{equation}\label{eq:operator49}
 Q_t(v;x,a)=c_t(x,a)+\beta_t\rho_t(v(F_t(x,a,Z))),\qquad
 \mathcal T_tv(x)=\inf_{a\in A_t(x)}Q_t(v;x,a).
\end{equation}
Throughout the deterministic comparison, the evaluations $\rho_t$ are monotone and cash invariant on a common domain containing all functions and constant translates used below. They are therefore nonexpansive in the supremum norm. Continuous objectives, effective compact feasible correspondences and measurable minimizers, or measurable arbitrarily near-minimizers with their error charged, are assumed. The stipulated conditional recursion gives Bellman comparison against all admissible adapted controls. In the earlier expected-cost experiments $\rho_t$ is conditional expectation, costs are bounded and innovations have their original continuous uniform law. Section~\ref{sec:continuous95} evaluates the same investment technology under an explicitly nested AVaR preference and retains the original continuous innovation law. The direct-cost inference in Section~\ref{sec:decisions49} is specific to expectation; it is not automatically an inference theorem for nonlinear recursive utility or games.

A Neural Bellman Operator is a construction and evaluation procedure for \eqref{eq:operator49}. It forms targets using its \emph{own} fitted $f_{t+1}$, returns a neural continuation and a feasible policy, and records an economic error allowance. It does not receive the optimal continuation as a training label. Arbitrary trained networks can be evaluated by the first result below. The subsequent constructive backend specifies how a network and its actor are obtained without assuming successful nonconvex optimization.

\subsection{Action-dependent errors and complete-policy loss}
Let $H(a)$ be a true current cost and $\widehat H(a)$ its continuation approximation. If $\widehat a$ is $\delta$-optimal for $\widehat H$ on the same feasible action set, then
\begin{equation}\label{eq:center49}
 H(\widehat a)-\inf_aH(a)\le\delta+
 \operatorname{osc}_{a}\{\widehat H(a)-H(a)\}.
\end{equation}
Indeed, compare $\widehat a$ with an arbitrarily near-optimal true action and cancel the approximate objectives. The infimum limit gives the result. Additive continuation levels cancel; action-dependent discrepancies do not. Under a change in future policies, the continuation must be transported to the new economic problem before this inequality licenses reuse. The full development edition retains the original quantitative transport, recursive-domain and economic-application results. No transport bound is presumed valid after an unaccounted change of constraints or preferences.

\begin{proposition}[One-sided policy sandwich]\label{prop:sandwich49}
Suppose, uniformly in state,
\[
 f_t-\mathcal T_tf_{t+1}\le u_t,\qquad
 \mathcal T_t^\pi f_{t+1}-f_t\le v_t,
 \qquad \ell_T\le f_T-g\le u_T.
\]
Then the actual feasible policy satisfies
\begin{equation}\label{eq:sandwich49}
 0\le J_0^\pi-V_0^*\le
 \sum_{t<T}B_t(u_t+v_t)+B_T(u_T-\ell_T).
\end{equation}
The bound is invariant to date-specific constant shifts of the fitted continuations, provided the associated residual accounts are shifted consistently.
\end{proposition}
\noindent The proof is in Supplement, Section~\ref{supp:proof-core49-1}.


A sharper valid residual account for the same continuation and the same deployed actor sharpens knowledge of the policy. It does not retrain that policy or change its realized cost. Conversely, subtracting two upper bounds in \eqref{eq:sandwich49} does not identify the difference in actual costs. These distinctions govern both the historical comparisons and the new experiments.

\section{Constructing and compiling a feasible neural policy}\label{sec:construction49}
Assume the primitive objective has, for feasible state--action pairs, state modulus $A_t^Q=C_{x,t}+\beta_tF_{x,t}L_{t+1}$ and action modulus $D_t=C_{a,t}+\beta_tF_{a,t}L_{t+1}$. Let
\[
 d_H(A_t(x),A_t(z))\le K_t^A\|x-z\|_1,
 \qquad L_t=A_t^Q+D_tK_t^A,\qquad L_T=L_g.
\]
The Hausdorff condition and the primitive moduli imply that $\mathcal T_tf_{t+1}$ is $L_t$-Lipschitz whenever $f_{t+1}$ is $L_{t+1}$-Lipschitz. All moduli and numerical integration procedures are supplied by the economic primitives, not fitted from successful outcomes.

At state nodes $x_{t,i}$ with covering radius $h_t$, form effective feasible $k_t$-nets. Evaluate their own-future objectives with certified error at most $e_t$. Minimize the computed labels and retain the corresponding \emph{original} action $a_{t,i}$. A Borel repair $\Gamma_{t,i}(x)$ is feasible at $x$ and satisfies
\begin{equation}\label{eq:repair49}
 \|\Gamma_{t,i}(x)-a_{t,i}\|_1\le K_t^A\|x-x_{t,i}\|_1+\zeta_t.
\end{equation}
The allowance $\zeta_t$ concerns displacement; feasibility itself may not fail. Define
\begin{equation}\label{eq:cone49}
 f_t(x)=\min_i\{y_{t,i}+L_t\|x-x_{t,i}\|_1\},\quad
 i_t(x)=\min\arg\min_i\{y_{t,i}+L_t\|x-x_{t,i}\|_1\}.
\end{equation}
Return $\pi_t(x)=\Gamma_{t,i_t(x)}(x)$. Terminal labels approximate $g$ with error $e_T$ and use the same cone construction with modulus $L_g$.

\begin{theorem}[Own-future construction with feasible witnesses]\label{thm:construct49}
The functions in \eqref{eq:cone49} have exact affine--ReLU realizations and Lipschitz moduli $L_t$ independent of label error. The returned policy is feasible and measurable. It satisfies
\begin{align}\label{eq:enclosures49}
 -e_t\le f_t-\mathcal T_tf_{t+1}&\le D_tk_t+e_t+2L_th_t,\\
 \mathcal T_t^\pi f_{t+1}-f_t&\le e_t+D_t\zeta_t,\notag
\end{align}
and hence
\begin{equation}\label{eq:constructionloss49}
 J_0^\pi-V_0^*\le
 \sum_{t<T}B_t(D_tk_t+2e_t+2L_th_t+D_t\zeta_t)
 +B_T(2e_T+2L_gh_T).
\end{equation}
With effective covers, repairs and certified queries at the prescribed accuracy, allocating each nonnegative term a share of a positive economic tolerance gives a terminating finite construction. A fixed numerical precision or a capped finite experimental ladder need not contain that allocation.
\end{theorem}
\noindent The proof is in Supplement, Section~\ref{supp:proof-core49-2}.


The Lipschitz-extension principle is classical \citep{mcshane1934}. A native min-plus envelope and its exact ReLU circuit, using the same labels, witnesses and ties, are \emph{one method with two encodings}. They return the same function and policy. Neither Theorem~\ref{thm:construct49} nor a faster representation proves a uniquely neural capability or a training-success probability for Adam or L-BFGS. The mathematical increment is the complete feasible, own-future construction and implementation account; the empirical question is its work and policy value relative to conventional generators.


\paragraph{Primitive accuracy and resource growth.}
The complete primitive investment construction, its original-optimum accuracy theorem, and the dated horizon allocation are retained in Appendix Sections~\ref{sec:accuracy62} and~\ref{sec:complexity86}. Their state, action, innovation, and bit-growth costs remain explicit. The continuous recursive-risk investment analysis below supplies the corresponding original-law qualification used in the current comparison.


\section{Reference-centered risk kernels and complete-policy qualification}\label{sec:centeredrisk92}
A datewise worst-case error account can charge incompatible states at successive dates. Under nonlinear recursive preferences, replacing that account by an expectation under the physical transition law is not valid. We instead retain the conditional risk evaluation and center its dual representation at the actual stored NBO continuation. The resulting state-dependent certificate has an original-optimum interpretation and a finite occupation representation. It extends, rather than replaces, the signed qualification of Section~\ref{sec:risk91}.

\subsection{The centered operator}
Retain the original augmented economic state, feasible actions, nonnegative discounts, and bounded complete-policy recursions. At each $(t,x,a)$ suppose the original conditional evaluation admits the representation
\begin{equation}\label{eq:dualrisk92}
 \rho_{t,x,a}(v)=\sup_{q\in\mathcal Q_{t,x,a}}
 \left\{\int v(y)q(dy)-\psi_{t,x,a}(q)\right\}.
\end{equation}
Every $q$ is a probability measure on the original successor state space. The representation is assumed on a common bounded-function domain containing all continuations below, with the measurability and selection conditions needed for the original Bellman recursions. This is a hypothesis about the original preference primitive, not permission to replace a static risk objective by a nested one. Conditional risk duality and risk-averse dynamic programming are established principles \citep{ruszczynski2010risk,ruszczynski2014erratum}. For a nonconvex or non-cash-invariant aggregator lacking \eqref{eq:dualrisk92}, the earlier order and modulus results remain applicable.

Fix one qualified parameter $\theta$ and suppress its subscript. Its entire sequence of stored continuations is $f=(f_t)_{t=0}^T$, and its supplied complete feasible policy is $\pi$. Define, using the same $f_{t+1}$ in every term,
\begin{align}\label{eq:centered-def92}
 r_t(x,a)&=c_t(x,a)+\beta_t\rho_{t,x,a}(f_{t+1})-f_t(x),\notag\\
 k_t^f(x,a,q)&=\psi_{t,x,a}(q)+\rho_{t,x,a}(f_{t+1})
                         -\int f_{t+1}\,dq,\notag\\
 \mathcal E_{t,x,a}^f(z)&=\sup_{q\in\mathcal Q_{t,x,a}}
                   \left\{\int z\,dq-k_t^f(x,a,q)\right\}.
\end{align}
The penalty is nonnegative and has infimum zero, whether or not the supremum in \eqref{eq:dualrisk92} is attained. Direct substitution gives
\begin{equation}\label{eq:exact-center92}
 \mathcal E_{t,x,a}^f(z)
 =\rho_{t,x,a}(f_{t+1}+z)-\rho_{t,x,a}(f_{t+1}).
\end{equation}
In particular, $\mathcal E^f(0)=0$, and $\mathcal E^f$ is monotone, cash invariant, and nonexpansive in the uniform norm. Discarding $k^f$ is an upper relaxation, not this identity. Freezing one risk-maximizing kernel at $f$ is in general not even an upper relaxation.

\begin{theorem}[State-dependent original-optimum certificate]\label{thm:centeredrisk92}
Suppose verified functions satisfy $r_t^-(x,a)\le r_t(x,a)$ for every original feasible action and $r_t(x,\pi_t(x))\le r_t^{\pi,+}(x)$. Let bounded measurable $L_t,U_t$ satisfy
\begin{align}\label{eq:potential-cert92}
 L_T&\le g-f_T\le U_T,\notag\\
 L_t(x)&\le\inf_{a\in A_t(x)}
       \{r_t^-(x,a)+\beta_t\mathcal E_{t,x,a}^f(L_{t+1})\},\notag\\
 U_t(x)&\ge r_t^{\pi,+}(x)
                +\beta_t\mathcal E_{t,x,\pi_t(x)}^f(U_{t+1}).
\end{align}
Then, at every initial state,
\begin{equation}\label{eq:potential-loss92}
 f_0+L_0\le V_0^*\le V_0^\pi\le f_0+U_0,
 \qquad 0\le V_0^\pi-V_0^*\le U_0-L_0.
\end{equation}
The same statements hold at every intervening date and state. They are uniform over a qualified parameter class whenever all premises and the reported bound hold for each member with its own common continuation sequence and complete policy. The certificate permits endogenous states and actions changing at every date; it does not require a finite state space or independent deployment contexts.

With exact residuals and terminal discrepancies, equality recursions in \eqref{eq:potential-cert92} give $L_t=V_t^*-f_t$ and $U_t=V_t^\pi-f_t$. With scalar bounds from \eqref{eq:signed-input91}, their corresponding constant lower and upper sequences are feasible certificates here. Thus the exact state-dependent recursion is no looser than the signed scalar account. A strictly smaller bound can result without changing a fit, a feasible action, or an economic observation.
\end{theorem}

The last assertion concerns the best exact recursion or a verified improvement upon the constant certificates. It is not a claim that every approximately computed $L,U$ improves the scalar account. In implementation, the minimum of valid policy supersolutions and the maximum of valid optimum subsolutions retain validity, so an inferior new bound need not displace the existing one.

The lower inequality requires every original action. A discrete action search without a verified original-menu lower bound cannot establish it. For a fixed action, a feasible risk kernel supplies a lower bound on $\mathcal E^f(L)$; an upper bound requires control of the entire risk envelope. Numerical integration, feasible repairs, residual errors and penalty errors belong in these inequalities. If penalty enclosures are used, lower penalty endpoints enter an upper certificate and upper endpoints enter a lower witness. Approximate dual optimizers are accepted only through these independently checked inequalities.


\section{Continuous-law qualification in the investment model}\label{sec:continuous95}
The risk theorem becomes a numerical method only when its inequalities can be constructed from economic primitives. We now supply this construction for the original investment technology with a nested tail-risk preference. States and innovations remain continuous, every feasible action is covered, and the implemented terminal controller rather than a hypothetical exact optimizer determines policy cost. This application joins continuation construction, the original-optimum comparison and a representation comparison on one problem. It does not replace the general NBO construction or the other economic applications.

\subsection{Primitives and the object being approximated}
Let $d\geq2$ be even, $X=[0,1]^d$, and interpret subscripts cyclically. At dates $t=0,1$, the feasible investment is $0\leq a\leq A(x)=1/8+\bar x/8$, where $\bar x=d^{-1}\sum_jx_j$. Write $b_j=1/2$ for odd $j$ and $b_j=1/4$ for even $j$, and $s_j=(-1)^{j+1}$. The original transition is
\begin{equation}\label{eq:investment95}
 F_j(x,a,z)=\frac1{16}+\frac{x_j}{2}+\frac{x_{j+1}}8
       +\frac{x_j(1-x_{j+1})}{16}+b_ja+s_jz,
 \qquad z\sim\operatorname{Unif}[-1/32,1/32].
\end{equation}
Define
\begin{align}\label{eq:cost95}
 S(x)&=\frac2d\sum_j(x_j-5/8)^2+
       \frac1{4d}\sum_j(x_j-x_{j+1})^2+2(1/2-2\bar x)_+^2,\\
 g(x)&=S(x)+\frac2d\sum_j(x_j-5/8)^2,\qquad h(a)=a^2+4a^4.\nonumber
\end{align}
We use $\beta=15/16$ and, separately, $\alpha\in\{1/2,3/4\}$. Average value at risk is the upper-tail cost functional
\[
 \rho_\alpha(v)=\sup\left\{\int vq\,dP:
             0\leq q\leq(1-\alpha)^{-1},\ \int q\,dP=1\right\}.
\]
Here $P$ is the original uniform innovation law. The objective is recursively defined by $V_2=g$ and
\[
 V_t(x)=S(x)+\inf_{0\leq a\leq A(x)}
       \{h(a)+\beta\rho_\alpha(V_{t+1}(F(x,a,\cdot)))\}.
\]
It is not the static tail risk of total realized cost. A complete policy value replaces each infimum by the action prescribed by that policy, with the same recursive risk primitive.

\paragraph{Assumptions for the continuous construction.}
For this subsection, the primitives are exactly \eqref{eq:investment95}--\eqref{eq:cost95}; $\alpha$ is fixed before construction; the fitted continuation $f$ is a finite convex function on $X$; and every numerical lower or upper value is an outward enclosure. Feasible tail densities and feasible actions are retained explicitly. Terminal selection is the measurable, dyadically acquired controller described below, with its fallback included in the error bound. These concrete conditions imply compact feasible graphs, bounded costs, measurable minimizers, and the conditional-risk dual used in Theorem~\ref{thm:centeredrisk92}. They do not assert that every recursive preference or every trained network has this convex structure.

\begin{proposition}[Primitive curvature and state supports]\label{prop:curvature95}
For either risk level and every even $d$, the one-date objective
$Q_1(x,a)=S(x)+h(a)+\beta\rho_\alpha(g(F(x,a,\cdot)))$
is jointly convex on $X\times[0,1/4]$, with curvature bounded below by
\begin{equation}\label{eq:curvature95}
 \frac{4-5\beta/4}{d}\|v\|_2^2+2w^2
\end{equation}
in a direction $(v,w)$. Its value $V_1$ is convex on $X$ and satisfies
\begin{equation}\label{eq:lipschitz95}
 |V_1(x)-V_1(y)|\leq\frac{15}{d}\|x-y\|_1.
\end{equation}
Let $W$ be the objective evaluated under any fixed feasible risk density, at an anchor $(x_0,a_0)$, and let $(p_x,p_a)$ be its gradient there. For any $\lambda,\nu\geq0$, set $e=p_a+\lambda-\nu$. Then the following is a global affine minorant of $V_1$:
\begin{align}\label{eq:affine95}
 \ell(y)={}&W+\lambda(a_0-A(x_0))-\nu a_0-\frac{e^2}{4}\nonumber\\
 &+\left(p_x-\frac{\lambda}{8d}\boldsymbol{1}\right)^{\!\top}(y-x_0).
\end{align}
Neither optimal multipliers nor an exact optimizing action are required.
\end{proposition}

The curvature is obtained from the original nonlinear technology, not imposed on an interpolated value grid. The negative second-order term from $x_j(1-x_{j+1})$ is bounded using $|\partial_jg|\leq10/d$; the stage curvature dominates it. The allowance $e^2/4$ pays for imperfect stationarity. The proof in the Technical Supplement also establishes the global state modulus in \eqref{eq:lipschitz95}, including the state-dependent action constraint.

\subsection{An all-action certificate under the original innovation law}
Partition the innovation interval into $n$ equal bins $B_i$, with midpoints $z_i$ and probabilities $p_i=1/n$. An innovation bin is a set over which an inequality holds, not a replacement point mass in the economic model. Choose nonnegative $q_i\leq p_i/(1-\alpha)$ summing to one. These masses define a feasible density constant on each bin. For a fixed initial state and feasible proposal $a_0$, let $W_0$ be a lower bound on the objective under this density, obtained from globally valid affine supports of the continuation at the successor midpoints. Let $p$ enclose the derivative of the same supporting objective at $a_0$. Then
\begin{equation}\label{eq:actionlower95}
 \underline V_0(x)=\inf_{0\leq u\leq A(x)}
         \{W_0+p(u-a_0)+(u-a_0)^2\}
\end{equation}
is a lower bound on the original Bellman optimum when the supported continuation is $V_1$. The scalar quadratic is minimized over the full original interval. Interval uncertainty in $p$ is covered by taking the smaller minimum at its two endpoints, since the minimum is concave in $p$.

For a convex fitted continuation, midpoint supports give the lower risk witness and the larger endpoint value bounds the continuation throughout each bin. For $V_1$, Proposition~\ref{prop:curvature95} supplies midpoint supports and the allowance $15/(32n)$ covers each half-bin. Upper tail risk is the average of the largest $n(1-\alpha)$ upper bin bounds. This average is the exact risk of a dominating step function, hence an upper bound under the continuous law.

\begin{theorem}[Constructed original-optimum qualification]\label{thm:continuous95}
Suppose the supports and bin bounds above are outwardly verified. Suppose further that the complete terminal controller satisfies
$V_1\leq V_1^{\widehat\pi}\leq V_1+\varepsilon_1$ everywhere on $X$.
For every queried initial state, the lower value \eqref{eq:actionlower95} and the binwise upper bound $\overline J_0^{\widehat\pi}$ satisfy
\begin{equation}\label{eq:originalgap95}
 0\leq J_0^{\widehat\pi}(x)-V_0(x)
       \leq\overline J_0^{\widehat\pi}(x)-\underline V_0(x).
\end{equation}
For any stored convex continuation $f$, the same lower and upper constructions are admissible bounds for the centered-risk comparison of Theorem~\ref{thm:centeredrisk92}. Their validity does not depend on the representation of $f$, on a successful fit, or on the proposal being an exact minimizer.
\end{theorem}

To see the connection, write a verified continuation bracket as $\ell\leq V_1\leq u$ and use errors $L=\ell-f$, $U=u-f$. The centered identity gives
\[
 \rho_\alpha(f)+\mathcal E^f(L)=\rho_\alpha(\ell),\qquad
 \rho_\alpha(f)+\mathcal E^f(U)=\rho_\alpha(u).
\]
Joint evaluation avoids subtracting and then adding two separately rounded copies of $\rho_\alpha(f)$. The bound \eqref{eq:actionlower95} uses one feasible risk density to minorize the required all-action infimum; the policy upper bound controls the full risk envelope. Thus the application constructs sufficient inequalities for the original theorem. It does not claim that a non-neural continuation would receive a different validity theorem, or that computing the original-value witness is free.

The terminal controller observes each state coordinate rounded down to $40$ binary places. It accepts its floating proposal only after an outward all-action gap of at most $2^{-21}$ at that observed state; otherwise it uses rational bisection. The original-state transfer is less than $30\,2^{-40}$. Consequently $\varepsilon_1=2^{-20}$ covers the complete controller, including observation and fallback. Policy evaluation encloses this controller at every possible successor, not only at the midpoint states used by the integrator.


\section{Action-aligned acquisition and complete-policy work}\label{sec:action99}
A continuation is acquired to make a decision, not to minimize an approximation statistic in isolation. The retained investment experiment makes the distinction operational: every certificate-oriented fit reduces its validated residual span, whereas aggregate first-proposal rejections increase. We now construct an acquisition target in the units of the admission inequality. The construction uses the same original action set and risk law as the NBO policy account.

\subsection{A computable directional certificate}
Fix a state and date, and suppress their indices. Let $A\subset\mathbb R^k$ be nonempty, compact and convex, let $Q$ be the original action objective, and let $a\in A$ be the actually realized proposal. A verified supporting calculation supplies $L,U\in\mathbb R$, $s\in\mathbb R^k$ and $\mu>0$ such that
\begin{equation}\label{eq:support99}
 Q(a)\le U,\qquad
 Q(b)\ge L+s^{\mathsf T}(b-a)+\frac\mu2\|b-a\|_2^2
 \quad\text{for every }b\in A.
\end{equation}
These inequalities need not identify $s$ as a derivative of $Q$. In particular, an imperfect risk witness or continuation minorant can lower $L$ while retaining a useful supporting slope. Put $w=U-L$ and
\begin{equation}\label{eq:gamma99}
 \Gamma(a)=w-\min_{b\in A}
       \left\{s^{\mathsf T}(b-a)+\frac\mu2\|b-a\|_2^2\right\}.
\end{equation}
A certified lower bound on the displayed minimum gives a valid, possibly larger, $\Gamma$. An unverified optimizer status does not.

\begin{theorem}[Directional certificate for an NBO proposal]\label{thm:action99}
Under \eqref{eq:support99},
\[
 0\le Q(a)-\inf_{b\in A}Q(b)\le\Gamma(a).
\]
Let $\widehat s$ be a supporting slope of the surrogate used to propose $a$, let $N_A(a)=\{n:n^{\mathsf T}(b-a)\le0\ (b\in A)\}$, and suppose
$\|s-\widehat s\|_2\le d$. If a supplied $n\in N_A(a)$ satisfies
$\|\widehat s+n\|_2\le r$, then
\begin{equation}\label{eq:direction99}
 \Gamma(a)\le w+\frac{(d+r)^2}{2\mu}.
\end{equation}
For $A=[0,A_0]$, the exact minimum in \eqref{eq:gamma99} is attained at
$b=\operatorname{proj}_{[0,A_0]}(a-s/\mu)$.
\end{theorem}

The proof, including boundary actions, is in Supplement~\ref{supp:action99}. Equation~\eqref{eq:direction99} separates three trainable or computable quantities: the original-objective enclosure width $w$, the directional mismatch $d$, and the surrogate stationarity residual $r$. The residual span of the continuation is not one of these quantities. A readout can improve that span and worsen the supporting slope at a nearly binding action. Conversely, constant changes in a continuation need not affect any of them after consistent centering.

The supporting calculation is available under the original recursive-risk primitives. Suppose $F(b,\omega)=F(a,\omega)+G(\omega)(b-a)$, the stage cost is $\mu$-strongly convex, and the risk dual has an action-independent feasible witness $q$ on the original innovation space. If
\[
 V(F(b,\omega))\ge \ell(\omega)+v(\omega)^{\mathsf T}G(\omega)(b-a),
\]
then \eqref{eq:support99} holds with
\begin{equation}\label{eq:lift99}
 L=c(a)+\beta\left\{\int\ell\,dq-\psi(q)\right\},\qquad
 s=\nabla c(a)+\beta\int G^{\mathsf T}v\,dq,
\end{equation}
and any valid upper bound $U$ for the original risk objective. Integration and penalty errors must be enclosed in these quantities. The witness must remain feasible for every action being minorized; a witness feasible only at $a$ is insufficient when the risk set itself depends on the action.

In the investment economy, $G=(1/2,1/4)^{\mathsf T}$ at $d=2$, the innovation law and AVaR cap are action independent, and $h''\ge2$. Thus $\mu=2$ gives exactly the all-action quadratic of \eqref{eq:actionlower95}. The added result is the explicit link from an acquired directional support to its admission gap, rather than a new economic model or an assumption that an acquired value support is an exact derivative.


\section{An active-region compiler for the original Bellman oracle}\label{sec:compiler109}
A learned continuation must compete with the structure available in the original economy. In \eqref{eq:investment95}--\eqref{eq:cost95}, that structure includes an algebraic last-date Bellman oracle. Giving this oracle only to a neural method would confound representation with information. We derive it explicitly and give its implementation to every comparator.

Put $R=1/32$, $t=1-\alpha$, $r_j=F_j(x,0,0)$, and
\begin{align*}
 H_x&=\tfrac12-2\bar r,& k_x&=\frac{10}{d}\sum_js_jr_j,\\
 C_x&=\frac8d\sum_jb_j(r_j-5/8)
 +\frac1{2d}\sum_j(b_j-b_{j+1})(r_j-r_{j+1}).
\end{align*}
The following formula uses the original continuous law, not a fitted transition or a finite-shock replacement.
\begin{theorem}[Exact tail algebra and constrained Bellman regions]\label{thm:compiler109}
For $z\sim\operatorname{Unif}[-R,R]$ and $0<\alpha<1$, define
\begin{equation}\label{eq:tailpoly109}
 \mathcal R_\alpha(k)=\operatorname{AVaR}_\alpha(5z^2+kz)
 =\begin{cases}
 \dfrac{5R^2(1+\alpha+\alpha^2)}3+\dfrac{\alpha k^2}{20t},& |k|\le10Rt,\\[2mm]
 \dfrac{5R^2(4\alpha^2-2\alpha+1)}3+\alpha R|k|,& |k|\ge10Rt.
 \end{cases}
\end{equation}
The two expressions and their first derivatives agree at their junctions. For every even $d$ and every original state $x$, the last-date action derivative is
\begin{equation}\label{eq:cubic109}
\begin{split}
 \partial_a Q_1(x,a)={}&16a^3+2a+\beta\{C_x+41a/32-3(H_x-3a/4)_+\}\\
 &+\frac{5\beta}{4}\operatorname{clip}\!\left(\frac{\alpha(k_x+5a/4)}{10t},-\alpha R,\alpha R\right).
\end{split}
\end{equation}
It is continuous and strictly increasing, with slope at least two wherever its second derivative is defined. The interval $[0,A(x)]$ has at most four regions separated by
$H_x/(3/4)$ and the two solutions of $k_x+5a/4=\pm10Rt$.
On each region \eqref{eq:cubic109} is $16a^3+pa+q$, with $p\ge2$.
There is a unique constrained minimizer: zero or capacity when the corresponding endpoint derivative has the appropriate sign, and otherwise the unique interior root. With $u=\sqrt{p/48}$, that root on its region is
\[
 a=-2u\sinh\!\left\{\tfrac13\operatorname{arsinh}(q/(32u^3))\right\}.
\]
Coefficient preparation costs $O(d)$ arithmetic; region identification and the number of scalar inversions are bounded independently of $d$.
\end{theorem}

The complexity statement counts real arithmetic and scalar inversions, not unit-cost exact transcendental evaluation. The implementation uses floating inversion only to propose a dyadic action. It then applies the unchanged outward original-law supporting certificate, including the state-capacity multiplier and its stationarity allowance. A failed compiled certificate pays for its failed attempt and the original terminal oracle. Consequently compiler error can change work but cannot acquire certificate authority. The full terminal policy used in nested policy-cost evaluation is not replaced by the compiler. In particular, more accurate optimal-continuation information and a different deployed future policy are distinct interventions.

\paragraph{Gradients and the computational graph.}
The last-date value is continuously differentiable on the relative interior of the state cube, with the one-sided boundary interpretation, and
\begin{equation}\label{eq:envelope109}
 \nabla V_1(x)=\nabla_xQ_1(x,a^*(x))-\lambda(x)\nabla A(x),\qquad
 \lambda(x)=\boldsymbol1_{\{a^*=A(x)\}}[-\partial_aQ_1(x,a^*)]_+.
\end{equation}
Within a smooth interior action region,
$\nabla a^*=-\partial_{ax}Q_1/\partial_{aa}Q_1$; within a strictly active endpoint region it is respectively zero or $\nabla A$. No global differentiability of the actor across region changes is asserted. In particular, differentiating the rounded actor or treating a feasible witness slope as this exact gradient is incorrect. Supplement Section~\ref{supp:compiler109} proves the formulas, uniqueness, boundary rules and certificate invariance. The original general own-future construction remains Theorem~\ref{thm:construct49}; the new experiment isolates its acquired last-continuation slice, not a newly trained long-horizon network.

\section{Tail-aware action accuracy}\label{sec:tail109}
A critic can match action derivatives pointwise in the innovation and still propose the wrong risk-averse action. The missing quantity is the change in the adverse tail. This section refines the directional account in Theorem~\ref{thm:action99}; it does not replace an original-law witness by a training statistic.

Fix a state and a realized action $a$. On the same original probability space let $Z_a$ be the true continuation outcome and $H_a$ an acquired surrogate. In a neighborhood of $a$, suppose both are almost surely differentiable with integrably dominated derivatives. At $a$ assume their $\alpha$-quantiles are unique, carry no mass, and cut off upper tails of probability $t=1-\alpha$, where $0<\alpha<1$. These are local assumptions, not automatic properties of every fitted ReLU. Choose a constant $c$ and certified numbers $\epsilon,\delta\ge0$ such that
\[
 \|Z_a-H_a-c\|_\infty\le\epsilon,\qquad
 \|\nabla_aZ_a-\nabla_aH_a\|_\infty\le\delta.
\]
The constant is an anchor-level centering; no derivative of a data-dependent centering rule is taken. Write $\eta$ for the true quantile, $B_\epsilon=\{|Z_a-\eta|\le2\epsilon\}$, and, for an arbitrary vector $v_0$,
\[
 G_\epsilon=\operatorname*{ess\,sup}_{B_\epsilon}\|\nabla_aH_a-v_0\|_2,
 \qquad m_\epsilon=\min\{P(B_\epsilon),2\min(\alpha,t)\}.
\]
Set $G_\epsilon=0$ when $P(B_\epsilon)=0$.
\begin{theorem}[Centered tail-switch bound]\label{thm:tail109}
Under these hypotheses, whenever $G_\epsilon<\infty$,
\begin{equation}\label{eq:tailgradient109}
 \left\|\nabla_a\operatorname{AVaR}_\alpha(Z_a)
       -\nabla_a\operatorname{AVaR}_\alpha(H_a)\right\|_2
 \le \delta+G_\epsilon m_\epsilon/t.
\end{equation}
If $P(|Z_a-\eta|\le r)\le Cr^\kappa$ for the relevant $r$, the tail term is at most $G_\epsilon C(2\epsilon)^\kappa/t$. If additionally
$\|\nabla_aH_a-v_0\|_2\le L|Z_a-\eta|$ on $B_\epsilon$, it is at most
$LC(2\epsilon)^{\kappa+1}/t$.

Suppose the original current objective $Q$ is differentiable and $\mu$-strongly convex on a compact convex action set $A$, and its stage cost is unchanged in the surrogate $\widehat Q$. If $n\in N_A(a)$ and $\|\nabla\widehat Q(a)+n\|_2\le r_a$, then
\begin{equation}\label{eq:tailaction109}
 0\le Q(a)-\min_AQ
 \le \frac{\{r_a+\beta(\delta+G_\epsilon m_\epsilon/t)\}^2}{2\mu}.
\end{equation}
\end{theorem}

The proof localizes changed tail membership to a quantile band and uses the zero total mass of the difference between two maximizing densities. Thus a constant component of the surrogate action gradient cancels: tail \emph{gradient oscillation}, rather than its overall magnitude, controls the switching cost. Under a positive margin exponent, joint control of the centered value error and directional error gives a rate for action loss. A finite validation sample does not itself verify $C$, $\kappa$ or the essential suprema. At atoms or nonunique tails the stated differentiability theorem is unavailable; the feasible-density minorants of Theorem~\ref{thm:action99} remain the appropriate certification device. When $Z_a$ uses an optimal continuation, \eqref{eq:tailaction109} bounds the current-action gap; an imperfect implemented future policy requires its separate recursive allowance before interpreting the bound as complete-policy regret. Numerical witness slack is paid there through $U-L$, not silently dropped from \eqref{eq:tailaction109}.

\paragraph{An exact diagnostic.}
Let $z$ be uniform on $[-1,1]$, $\alpha=1/2$, $a\in[-1,1]$, and stage cost $a^2$. Take $Z_a=(a+1)z$ and $H_a=(a-1)z$. Their action derivatives are identically $z$, yet the original and surrogate optimal actions are $-1/4$ and $1/4$. Using the surrogate action incurs original loss $1/4$. This continuous-law diagnostic, independently integrated in rational arithmetic, rules out any general action-loss bound that vanishes with directional error alone. It is not the investment economy, a new population sample, or an empirical neural advantage. Its conceptual provenance in the earlier tail-aware research protocols is retained.

\paragraph{Acquisition target and novelty.}
For a family satisfying verified margin assumptions, a useful acquisition target is the bound in \eqref{eq:tailaction109}, with envelope slack and actual qualification work added. The earlier prospective finite-menu selection theorem then applies to complete services with that verified certificate. Our executed comparison below does not retroactively become a training experiment for this new loss: its frozen acquisition remains least squares. CVaR gradient estimation and decision-focused learning already have established formulations \citep{Tamar2015,LuoDelage2026,ElmachtoubGrigas2022}. The present result is a centered, tail-localized perturbation bound tied to the same constrained Bellman action and a valid complete-policy account. The static trajectory-CVaR objective of \citet{LuoDelage2026} is not the nested AVaR preference used here; its empirical gains cannot be transferred between those objectives.

\section{Safeguarded policy improvement and resolution}\label{sec:improvement109}
Qualification and improvement answer different questions. A completed policy may satisfy a loss tolerance without minimizing cost, and a rejected early proposal is not the completed conventional baseline. We retain a qualified incumbent before attempting any improvement.

\begin{proposition}[Completed-incumbent improvement]\label{prop:safe109}
Fix the original state, law, feasible action set and complete future policy. Suppose two feasible current actions have verified complete-policy cost intervals $[l,u]$ and $[l',u']$. Replacing the incumbent only when $u'<l$ guarantees a strict cost decrease of at least $l-u'>0$. Repeating this rule preserves weak improvement relative to the initial \emph{completed} incumbent. Requiring every replacement also to pass the original-optimum and policy-width targets preserves those targets. A refused trial leaves the incumbent unchanged but still incurs its full work.
\end{proposition}

This rule is an interval implementation of a descent safeguard, not a new claim that descent is uniquely neural. In the executed service, three projected support steps are attempted, with fractions $1,1/2,1/4$, after the original fallback and refinement have completed. An original-optimum lower bound can be reused across actions at the same state. A policy-cost interval cannot. The projection generates candidates; only the outward gate grants a policy change. Both learned and online services receive exactly this option.

\subsection{When local improvement becomes recursive improvement}
A finite catalogue does not establish a statewise policy theorem. To state the required condition, consider bounded measurable functions, bounded costs, $0<\beta<1$, and original fixed-policy operators
\[
 T_\pi v(x)=c(x,\pi(x))+\beta\rho_{x,\pi(x)}(v(F(x,\pi(x),Z))).
\]
Assume the operators preserve bounded measurability, the one-step risks are monotone and cash invariant, and the Bellman infimum admits measurable minimizers or effective near-minimizers with their errors charged. Assume also a common dual representation, at each fixed $(x,a)$,
$\rho_{x,a}(v)=\sup_{q\in\mathcal Q(x,a)}\{\int v\,dq-\psi_{x,a}(q)\}$,
where each $q$ is a probability law. The dual set may depend on the action, but the two continuations compared at that action use the same set and penalty.

\begin{theorem}[Recursive improvement with an economic gain account]\label{thm:improve109}
The original Bellman operator $T=\inf_\pi T_\pi$ and every $T_\pi$ are $\beta$-contractions. They have unique bounded fixed points $V^*$ and $J^\pi$. Suppose the \emph{uniform} improvement residual satisfies
\[
 0\le d(x)\le J^\pi(x)-T_{\pi'}J^\pi(x)\le e(x)
\]
for bounded nonnegative $d,e$. Let $P^-v=\inf_{q\in\mathcal Q(x,\pi'(x))}\int v(F)\,dq$ and let $P^+$ replace the infimum by a supremum; assume these maps preserve bounded measurability. Their unique discounted fixed points
\[
 D^-=d+\beta P^-D^-,\qquad D^+=e+\beta P^+D^+
\]
satisfy
\begin{equation}\label{eq:gain109}
 0\le D^-\le J^\pi-J^{\pi'}\le D^+.
\end{equation}
These gain accounts use risk-distorted future laws, not physical state frequencies.

If a sequence of such uniformly safe updates additionally satisfies
$T_{\pi_{k+1}}J^{\pi_k}\le TJ^{\pi_k}+\varepsilon_k$, then
\begin{equation}\label{eq:policyiteration109}
 \|J^{\pi_{k+1}}-V^*\|_\infty
 \le\beta\|J^{\pi_k}-V^*\|_\infty+\varepsilon_k.
\end{equation}
In particular $\varepsilon_k\to0$ implies convergence to the original optimum. A verified bound $0\le J^\pi-TJ^\pi\le(1-\beta)\varepsilon$ is a valid stopping certificate for $\|J^\pi-V^*\|_\infty\le\varepsilon$.
\end{theorem}

Risk-sensitive policy iteration is established theory \citep{ruszczynski2010risk,ruszczynski2014erratum}. Here \eqref{eq:gain109} states which original-objective residual an outward improvement gate must establish and how its economic gains propagate. The finite executed root catalogue verifies Proposition~\ref{prop:safe109}, not the uniform premises of Theorem~\ref{thm:improve109}. Neither three trial steps nor successful critic fitting establishes global policy-iteration convergence.

\subsection{Paid precision and a stopping rule}
A strict descent gate can be undecided because two policy intervals overlap. Suppose a fixed incumbent and candidate have actual improvement $\Delta=J-J'$, and refinement at resolution $n$ gives intervals containing their respective costs, with widths at most $A/n^2$ and $A'/n^2$. Then
\begin{equation}\label{eq:precisiongate109}
 l_n-u'_n\ge\Delta-(A+A')/n^2.
\end{equation}
For $\Delta>0$, doubling $n$ until $n>\sqrt{(A+A')/\Delta}$ certifies improvement. With $n_0=1$ and cost at most $cn$ per refinement round, the total paid work is less than
$4c\max\{1,\sqrt{(A+A')/\Delta}\}$, apart from fixed setup. If the proposed action has loss and width targets too, those must be verified separately. If $\Delta=0$, no finite strict-descent certificate is promised. Moreover, from interval information alone, overlap is an information obstruction: an algorithm cannot certify strict descent for every pair of costs consistent with two overlapping intervals. Thus the claimed resolution scale belongs to this interval interface, not a new universal lower bound for all economic oracles.

For an operational stopping decision, let $b\le V^*(x)$ be valid and let $u$ bound the incumbent cost. Its remaining possible cost reduction is at most $u-b$. If a known nonnegative compute price converts the unavoidable expense of a contemplated improvement operation to at least $u-b$ in the same units, that operation cannot provide strictly positive net gain at this single decision. For a reused policy, multiply the qualified gain by its declared workload before comparing with acquisition expense. This dominance rule is implementable once costs and prices are supplied; it is neither an empirical adoption claim nor a globally optimal search policy over unknown future discoveries. Supplement Section~\ref{supp:improve109} gives complete proofs.

\section{Certified remaining gain and paid stopping}\label{sec:budget110}
A completed NBO service should not keep paying for improvements that cannot be economically material. The relevant stopping quantity is not a critic's training residual. It is an upper bound on the remaining improvement in the original decision problem. We turn that bound into a statewise allocation rule, with an explicit price in possible foregone policy improvement. The rule is available to every representation; it is not a neural-specific optimization principle.

\subsection{A certificate interface in units of original cost}
Fix a state, the original information and law, and the same implemented future policy. An incumbent action is already feasible and qualified. Its cost is $J$, its verified interval is $[l,u]$, and $b\le V^*$ is a verified lower bound for the original optimal value. Let
\[
 e=\operatorname{up}(u-b),\qquad e\ge0,
\]
where $\operatorname{up}$ denotes a genuine outward arithmetic bound, not a favorable rounding of an optimizer's return value. Qualification still means the original loss and policy-width targets. Choose an additional, separately declared improvement budget $\theta\ge0$.

\begin{theorem}[Remaining-gain stopping and a complete-service comparison]\label{thm:budget110}
If $e\le\theta$, retaining the incumbent and stopping all further improvement at that state preserves its qualification and, for every feasible original policy $\nu$, gives
\begin{equation}\label{eq:stopgain110}
 (J-J^\nu)_+\le e\le\theta.
\end{equation}
The competitor may have a better future policy: the lower bound $b$ concerns the original optimum, not only the fixed-future action optimum.

Consider a finite safeguarded improvement procedure and a shortened version that starts from the same completed incumbent. The shortened procedure uses the same trial and acceptance rules until it stops by \eqref{eq:stopgain110}; at an un-stopped state suppose its complete trace agrees with the full procedure. Denote the returned policy costs by $J^{\rm short}$ and $J^{\rm full}$. Then
\begin{equation}\label{eq:short110}
 (J^{\rm short}(x)-J^{\rm full}(x))_+
 \le\boldsymbol1_{\{\text{stopped at }x\}}e_{\rm stop}(x).
\end{equation}
For a declared finite catalogue with nonnegative weights $w_i$ summing to one, the weighted mean positive-part loss is at most
$\sum_{i:\,\mathrm{stopped}}w_i e_{\rm stop}(x_i)$.
For identical returned actions and the same future policy, the cost contrast is exactly zero, not the width of two independently subtracted intervals.

From $[l,u]$ and $b$ alone, the ceiling $u-b$ cannot generally be reduced. This is sharpness of the declared certificate interface, not an information lower bound against a solver given additional formulas, derivative information or a sharper feasible-policy enclosure.
\end{theorem}

A final state can have $e>\theta$ after using all allowed trials without stopping. It must still meet the unchanged original qualification target. Thus $\theta$ is \emph{not} a relaxed or tightened substitute for the original policy-regret tolerance. The trace-agreement premise also matters in floating arithmetic: batching only the active rows can alter arithmetic order in an implementation that is not rowwise deterministic. Our protocol records and checks that issue; a mathematical stopping bound does not itself establish bitwise equality of all un-stopped computations.

The complete implementation first runs the compiled original-law service to qualification. It then considers the same three support-step fractions $1,1/2,1/4$ as the full safeguard. Before each round it stops a row when $e\le\theta$, retaining that row's action and certificate. Remaining rows pay for the original-law supporting query and fresh policy-cost enclosure. A changed action is accepted only by strict interval separation and both original qualification predicates. The stopping comparisons, active masks, refused trials and omitted trial slots are all recorded. Neither a cheap surrogate loss nor a noncertified estimate of improvement authorizes stopping.

\subsection{What propagates through a recursive problem}
A root-state guarantee with a fixed future policy is not a uniform multi-date policy guarantee. The following result states the additional condition. Use the bounded original discounted policy operators and the common dual risk sets of Theorem~\ref{thm:improve109}. Let $\pi_s$ be a shortened policy and $\pi_f$ a comparator. Suppose, \emph{uniformly over the relevant state domain},
\begin{equation}\label{eq:recursivebudget110}
 T_{\pi_s}J^{\pi_f}-J^{\pi_f}\le e,
 \qquad e\ge0.
\end{equation}
This is a residual evaluated against the comparator's continuation, not a training residual against a different critic. Define
\[
 P^+_{\pi_s}v(x)=\sup_{q\in\mathcal Q(x,\pi_s(x))}
                       \int v(F(x,\pi_s(x),z))\,dq(z).
\]
Assume this map preserves bounded measurability.
\begin{proposition}[Risk-distorted composition of foregone-gain budgets]\label{prop:compose110}
The unique bounded fixed point $E=e+\beta P^+_{\pi_s}E$ satisfies
\begin{equation}\label{eq:composedgain110}
 (J^{\pi_s}-J^{\pi_f})_+\le E.
\end{equation}
If $e\le\theta$ uniformly, then $E\le\theta/(1-\beta)$. For a finite horizon, the corresponding backward account is
$E_t=e_t+\beta_tP^+_{\pi_s,t}E_{t+1}$, with the declared terminal positive-part discrepancy as $E_T$. Uniform scalar budgets compose by the dated discounted sum, not by an unqualified physical-state average.
\end{proposition}

The proof uses monotonicity and the common-risk-envelope difference bound, without exchanging the control infimum and risk supremum. The executed two-date catalogue verifies the fixed-future statement in Theorem~\ref{thm:budget110}; it does not establish the uniform premises in \eqref{eq:recursivebudget110} for an untested longer-horizon controller.

\subsection{A computable economic trade-off, not an invented purchase price}
If an improvement operation has an unavoidable charge $c$ expressed in the same economic units as $J$, and $c\ge e$, no replacement can generate a strictly positive net gain at this decision. This is a dominance rule, not a globally optimal policy for learning which future computations are useful. When a continuation is reused for $B$ decisions, its acquisition cost and the appropriate joint or weighted gain bound must be compared in common units; $B$ is not supplied for free by a benchmark.

Our primary numerical budget is $\theta=2^{-16}$, with all three prespecified sensitivities $2^{-12},2^{-14},2^{-18}$ reported. It is an engineering foregone-gain budget, not an empirically calibrated welfare margin. The comparison charges all common acquisition and all active and refused work. It reports the actual cost sacrifice against the full safeguard, including positive sacrifice, alongside the computational saving. This directly answers a delivered-service question that a smaller global residual span alone cannot answer. Complete proofs follow in Supplement Section~\ref{supp:budget110}.

\paragraph{Recursive utility and the computational graph.}
Supplement Section~\ref{supp:bridge109} proves a Brownian Markov recursive-utility verification result, including well-posedness, the critic graph $z_u=\sigma^{\mathsf T}\nabla u$, and terminal or stopping-boundary errors. It is a delimited NDU/FBSDE research interface, not an equivalence between different economic objectives. The information statements below retain their declared oracle restrictions; formula or derivative access is not free point-value information.

\section{Decision-scale numerical and statistical precision}\label{sec:precision99}
The continuous qualification uses the original terminal-controller allowance and all-action lower witness, not just a small surrogate residual. Convex midpoint and endpoint risk rules supply a width of at most $KR/\{(1-\alpha)n^2\}$ with $2n+1$ continuation queries. Theorem~\ref{thm:nonuniform99} below extends this account beyond the uniform law. The full uniform-rule statement, numerical endpoint charges and tensor extension remain in Section~\ref{sec:precision96} of the Development Appendix.

For the frozen population comparison, every complete service is qualified to the original-optimum allowance $\tau=1/256$. The design uses 4,096 independent initial states at each risk level and twelve simultaneous contrasts, with numerical half-width and sampling allowance totaling $3/8192$. Theorem~\ref{thm:power96} states the exact simultaneous event and its independence requirements. Those population intervals belong to the frozen actors, not to subsequently refitted candidates. The acquisition comparison instead retains the deterministic contrast support $[-\tau,\tau]$.

For a workload of $B$ decisions, the purchase account combines the per-decision policy contrast, $B$ times that contrast, full-workload computation and one-time acquisition. Theorems~\ref{thm:purchase97} and Proposition~\ref{prop:purchase96} retain the complete interval and strict-boundary statements. Historical timing ranges are descriptive inputs to a conditional price exercise. The exact decision mechanism in Section~\ref{sec:decision99} will provide an identified model-specific policy contrast; it will not turn these historical clock ranges into confidence intervals or external prices into estimated welfare weights.

\section{Information resolution under nonuniform laws and randomization}\label{sec:information99}
The information required to qualify a continuation is distinct from the information supplied when its formula is already known. This section studies point-value access to a convex, $K$-Lipschitz continuation on $[-R,R]$, with $K,R>0$. The continuation formula, derivatives and exact Bellman labels are not supplied for free. The original uniform-law result, including its complete proof and constants, is retained as Theorem~\ref{thm:information97} in the Development Appendix. We extend the law and distinguish two nonequivalent meanings of randomized certification.

\subsection{A known nonuniform innovation law}
Let $P$ have known density $p$ satisfying $0<m\le p(z)\le M_p<\infty$ on $[-R,R]$. Fix $0\le\alpha<1$, put $t=1-\alpha$, let $q$ be the $\alpha$-quantile of $P$ (with $q=-R$ at $\alpha=0$), and set $L=R-q$. Let $\mathcal C_{K,R}$ denote the convex $K$-Lipschitz functions on this interval.

\begin{theorem}[Nonuniform-law information and enclosure]\label{thm:nonuniform99}
Every deterministic adaptive procedure using at most $N$ point values and reporting an interval valid for all consistent $f\in\mathcal C_{K,R}$ has an input with reported width at least
\begin{equation}\label{eq:lower99}
 \frac{mKL^2}{12t(N+1)^2}
 \ \ge\ \frac{mKt}{12M_p^2(N+1)^2}.
\end{equation}
For an upper bound, partition $[-R,R]$ into $n$ equal cells. Suppose the original cell masses and conditional cell means are available exactly. The mass-weighted conditional-mean rule and the mass-weighted endpoint secant rule bracket $\operatorname{AVaR}_\alpha^P(f)$ and have width at most
\begin{equation}\label{eq:upper99}
 \frac{2M_pKR^2}{t n^2}.
\end{equation}
They use at most $2n+1$ distinct continuation values. Thus the deterministic error-versus-information rate is $N^{-2}$ for each fixed admissible law and fixed $\alpha$. Certified, rather than exact, law moments require their additional numerical allowances.
\end{theorem}

The lower bound uses indistinguishable convex tails under the original density, not a probability transformation that might destroy convexity. The upper bound applies conditional Jensen and secant inequalities to $(f-\eta)_+$ before minimizing over the AVaR threshold $\eta$. The moments are law information and their acquisition cost is explicit. With $m=M_p=1/(2R)$ and $L=2Rt$, both constants reduce to their uniform-law counterparts. Their opposite behavior as $\alpha$ tends to one remains; no tail-uniform constant is identified.

\subsection{Sound random intervals and randomized point estimates}
A randomized algorithm can mean a zero-failure enclosure or a statistical approximation. We call an interval procedure \emph{pathwise sound} if, outside one null set of random tapes independent of the input, it contains the risk of every admissible input consistent with the tape's transcript. This is stronger than requiring coverage $1-\delta$ separately for each input.

\begin{proposition}[Randomized information with explicit quantifiers]\label{prop:random99}
The first lower bound in \eqref{eq:lower99} holds on one fixed quadratic-tail input for every sound random tape of a pathwise-sound interval procedure. It therefore also bounds expected width on that input.

For ordinary randomized point-valued procedures under the uniform law, let $A(f)$ use at most $N\ge1$ point values on every realization. Then
\begin{equation}\label{eq:randomlower99}
 \sup_{f\in\mathcal C_{K,R}}
 \left\{\mathbb E\left|A(f)-\operatorname{AVaR}_\alpha^P(f)\right|^2\right\}^{1/2}
 \ge\frac{KR(1-\alpha)}{960\,N^{5/2}}.
\end{equation}
The expectation is over the algorithm's randomization. This is a lower bound, not a claim that $N^{-5/2}$ is the optimal randomized rate on the whole class.
\end{proposition}

For the first assertion, the adverse quadratic is fixed before the random tape; its indistinguishable chord continuation can depend on the realized transcript because soundness is simultaneous over inputs. That reasoning cannot be applied to a merely high-probability interval. The second assertion instead uses a fixed finite distribution of convex perturbations and the conditional variance of unrevealed signs. Its proof supplies all constants and treats adaptive queries. Classical information-based complexity already uses indistinguishable inputs and randomized bump distributions \citep{Novak2015Integration97}; the result here specializes those tools to the convex monotone AVaR tail and states which certification quantifier survives randomization.

\subsection{Dated resolution and the Bellman resource account}
Suppose, on a fixed acquired state--action graph, date $j$ contributes a verified nonnegative allowance $a_j/n_j^2$ to the original-policy tolerance and incurs $c_j n_j$ continuation-query work, with $a_j,c_j>0$. The coefficients include discounting, the appropriate tail-risk factor and the number of original objectives to be evaluated. The continuous allocation problem has value
\begin{equation}\label{eq:allocation99}
 \min_{n_j>0:\,\sum_j a_j/n_j^2\le\varepsilon}
       \sum_j c_jn_j
 =\frac{\left(\sum_j a_j^{1/3}c_j^{2/3}\right)^{3/2}}{\sqrt\varepsilon}.
\end{equation}
Rounding each minimizing $n_j$ upward preserves the tolerance and adds less than $\sum_jc_j$ query work. This is a constructive composition of certified error accounts, not an end-to-end Bellman lower bound. Acquiring the state--action graph, the neural continuation, law moments, exact arithmetic and feasible owners remains additional work. With multiple innovations a direct tensor implementation still has product point counts. The original horizon and vector-control results in Proposition~\ref{prop:horizon95} and the Development Appendix are retained, including their adverse executed scaling evidence.

\section{A decision-theoretic value of acquired oblique geometry}\label{sec:decision99}
An approximation separation matters economically only when it changes an implemented decision. We give such a realization for the original oblique continuation, without altering the retained investment comparison. This is an exactly solved economic mechanism, not an empirical calibration or a timing experiment.

A decision maker observes $x=(x_1,x_2)\in[0,1]^2$ and a posted price $p\in[0,3]$. Purchasing one durable claim yields next-date benefit
\begin{equation}\label{eq:benefit99}
 v(x)=1+(x_1+x_2-1)_+.
\end{equation}
The feasible action is $b\in\{0,1\}$, and total cost is $b\{p-v(x)\}$. This is a two-date Bellman problem with a deterministic terminal payoff and unit discount. A stored continuation $h(x)\in[0,3]$ implements $b_h(x,p)=\boldsymbol1\{p\le h(x)\}$. The original optimum buys at $p\le v(x)$. A price tie has zero loss. Let states be uniform on the square and prices uniform on $[0,3]$, independently.

\begin{theorem}[Exact policy-loss separation]\label{thm:decision99}
For every measurable $h:[0,1]^2\to[0,3]$, the conditional and unconditional policy regrets are
\begin{equation}\label{eq:decisionloss99}
 \mathbb E_p[J^{b_h}(x,p)-J^*(x,p)]
       =\frac{(h(x)-v(x))^2}{6},\qquad
 \mathbb E[J^{b_h}-J^*]=\frac{\|h-v\|_{L^2}^2}{6}.
\end{equation}
Among arbitrary square-integrable additive continuations and among polynomials of total degree at most two, respectively, the optimal expected regrets are
\begin{equation}\label{eq:separation99}
 \frac1{540}\quad\text{and}\quad\frac1{5400}.
\end{equation}
They are attained by the admissible functions
\begin{align}
 h_A(x)&=\frac56+\frac{x_1^2+x_2^2}{2},\label{eq:additive99}\\
 h_Q(x)&=\frac{17}{15}-\frac35(x_1+x_2)
              +\frac12(x_1^2+x_2^2)+\frac65x_1x_2.\label{eq:quadratic99}
\end{align}
A single oblique ReLU with the affine constant in \eqref{eq:benefit99} attains zero regret. A conventional implementation of the identical oblique hinge also attains zero regret and the same decisions.
\end{theorem}

The shift by one is part of the economic construction. It keeps the projection optima inside the posted-price support, so clipping does not secretly enlarge a comparator class. The proof computes the additive conditional-expectation projection and the six-dimensional polynomial Gram system exactly. It establishes a policy comparison against the best member of each specified class, rather than comparing arbitrarily chosen fits. The scalar price integration is the step that turns continuation error into action loss; no such identity is presumed in the investment economy.

The sup-norm comparison also has a decision interpretation. Under the same conditional price law, the smallest worst-state expected regrets for the two classes are $1/96$ and $1/1536$, respectively. These follow from the sharp errors $1/4$ and $1/16$ in Proposition~\ref{prop:ridge97}, with its explicit feasible attaining functions shifted by one. The population optima in \eqref{eq:additive99}--\eqref{eq:quadratic99} are different from the minimax optima. An acquisition objective must specify which context criterion it is intended to serve.

\subsection{Acquisition, conventional geometry and purchase}
The mechanism identifies a value of the oblique continuation geometry, not a value of relabeling its formula as a network. A piecewise-affine conventional representation on the two triangles separated by $x_1+x_2=1$ is also exact. Convex regression with that hinge already in its dictionary contains the same exact continuation. Those controls are included in the exact comparison. Discovering the separating hyperplane from unknown data, fitting a larger dictionary or constructing an adaptive triangulation are different acquisition problems; their costs cannot be inferred from this formula-access calculation.

Suppose the exact oblique service has additional fixed acquisition cost $K\ge0$, incremental per-decision computation cost $c$, and fee $F\ge0$, all expressed in the units of this model's cost. Against the best quadratic continuation, the exact gain from $m$ independent decisions is
\begin{equation}\label{eq:exactpurchase99}
 G_Q(m)=m(1/5400-c)-K-F.
\end{equation}
When $c<1/5400$, the first strictly beneficial workload is
$\lfloor(K+F)/(1/5400-c)\rfloor+1$. This is a nonzero, exactly identified policy benefit with an explicit acquisition threshold. It is not an estimate of the price of computation in the retained investment application. Equal hinge and neural implementations share this gain when their complete costs agree.

\begin{table}[htbp]
\caption{Exact original-law policy loss in the oblique-continuation mechanism}\label{tab:decision99}
\centering\small
\begin{tabular}{@{}p{.52\linewidth}p{.17\linewidth}p{.22\linewidth}@{}}
\toprule
Continuation class or supplied geometry & Optimal mean regret & Information boundary\\\midrule
Arbitrary square-integrable additive functions & $1/540$ & Population projection\\
Total-degree-two polynomials & $1/5400$ & Six exact moments\\
One oblique ReLU and affine constant & $0$ & Supplied hinge\\
Identical conventional oblique hinge & $0$ & Same supplied hinge\\
Piecewise affine on the two separating triangles & $0$ & Same separating line\\
Convex dictionary containing the oblique hinge & $0$ & Contains exact function\\\bottomrule
\end{tabular}
\par\smallskip\parbox{.96\linewidth}{\footnotesize Notes: States are uniform on $[0,1]^2$ and independent prices uniform on $[0,3]$. Loss is evaluated against the original purchase optimum. The last four rows are equivalent exact geometric realizations, not four independent fitted experiments. The quadratic row uses six right-side moments and its known Gram matrix; moment acquisition is not free when the continuation formula is unknown. No timing, empirical calibration or dimension-wide neural dominance is inferred.}
\end{table}

The table is generated from two separately implemented exact-arithmetic calculations. The production path assembles monomial moments and solves normal equations by rational elimination. The independent path integrates polynomials after a triangular change of variables and uses an orthogonal shifted-Legendre basis. It imports neither the production integrator nor its solver. Both reproduce the fitted function, the original-law loss and the range checks. This independently checks acquisition for the stated finite-dimensional mechanism; it does not independently refit the R97 nonlinear dictionary search or complete the separately frozen R98 campaign.



\section{Original-policy costs and the price of refinement}\label{sec:study110}
\subsection{Economic identity, acquisition and prospective design}
The experiments use exactly \eqref{eq:investment95}--\eqref{eq:cost95}: two dates, two state coordinates, scalar state-dependent investment capacity, discount $15/16$, continuous uniform innovation and nested upper-cost AVaR at $\alpha=1/2$ and $3/4$. The actual terminal controller and its $2^{-20}$ uniform allowance are unchanged. Every returned policy must have original-optimum loss at most $1/256$ and a policy-cost interval width at most $1/8192$. Numerical integration encloses the continuous law rather than replacing it by a finite-shock economic model.

Every fitted arm receives the same original-model terminal-value enclosures at the 256 midpoints of a $16\times16$ grid. Their exact production interval midpoints are regression labels, and their acquisition is charged. The comparison is the acquired last-continuation slice of NBO; it does not newly train the general own-future construction at an unreported longer horizon. Evaluation uses all 64 midpoints of a disjoint $8\times8$ grid. Those catalogues were declared before execution, but the same design appeared in earlier research protocols. It is not an unseen statistical sample. The bounds below average the entire declared catalogue and are numerical intervals, not population confidence intervals. No R96 or R97 statistical interval is transferred to a newly fitted actor.

Rank-six quadratic, ReLU and coordinate-spline arms use their original 32-feature dictionaries, three nonlinear selections, nonnegative nonlinear coefficients, the declared numerical tie band and 30-bit rounding. A separately coded conventional hinge evaluates precisely the acquired ReLU geometry. Full degree-two and degree-three polynomials and bilinear interpolation broaden the conventional frontier; their unequal ranks and memory are recorded, not forced into an artificially weak neural comparison class. Polynomial and grid action proposals use the fixed scan and local scalar search. Their numerical optimization status is never original-law certificate authority.

Independent feature-matrix assembly, nonlinear active-set enumeration and pivoted QR reproduce all nonexcluded candidate fitting objectives. All six rank-six feature selections and rounded coefficient vectors agree exactly in this execution. Full polynomials are checked with a separate factorization. Every serialized model is reloaded and its predictions checked. Shared floating linear-algebra libraries are disclosed: this is independent algorithmic reconstruction, not exact optimization certification or independent human authorship. Hinge and ReLU share and pay the same mathematical acquisition and are not counted as independent learning successes.

\paragraph{Two attributable campaigns.}
The first campaign reproduces unchanged executable source frozen in R108 at \texttt{29180a4}, including the four original complete services: legacy, compiled, and a safeguarded version of each. Its current run has its own raw files, host records and clocks. The second campaign was prospectively deposited at \texttt{f43602d}, with actual executable bytes frozen at \texttt{5e1eae8}, before its primary execution. It uses those newly reproduced actor payloads without refitting. The new six-service comparison has plain compiled and full compiled-safe controls plus gain budgets $2^{-12},2^{-14},2^{-16},2^{-18}$; the primary budget is $2^{-16}$. Earlier aggregate findings informed the research question, a disclosed chronology rather than a claim of blind design. No other frozen experiment is completed or amended by either campaign.

\subsection{Complete decision services}
Every service first completes the original 128-bin qualification, paid online repair where needed, refinement ladder and exact fallback. The full safeguard then attempts all three original-support fractions $1,1/2,1/4$. A trial changes a policy only when its original-policy upper cost is strictly below the incumbent's lower cost and both original targets pass. The shortened service adds only the remaining-gain predicate of Theorem~\ref{thm:budget110}. Active rows still pay for the same original-law support and candidate certificate; stopped rows retain the completed incumbent and its recorded gain ceiling. All refused trials, masks and stopping comparisons are charged.

\begin{table}[htbp]\caption{Primary remaining-gain budget: complete work and actual foregone improvement}\label{tab:budget110}
\centering\scriptsize
\begin{tabular}{@{}llrrrrr@{}}\toprule
$\alpha$ & Arm & Full-safe & Budgeted & Omitted & Updates & Extra mean cost $\times10^6$\\\midrule
$1/2$ & Online & 0.4251 & 0.1486 & 192/192 & 0 & [0.000, 0.000]\\
$1/2$ & Quadratic (rank 6) & 0.3686 & 0.1687 & 133/192 & 44 & [0.454, 0.928]\\
$1/2$ & ReLU (rank 6) & 0.4074 & 0.2074 & 153/192 & 33 & [0.195, 0.334]\\
$1/2$ & Equivalent hinge & 0.4071 & 0.2049 & 153/192 & 33 & [0.195, 0.334]\\
$1/2$ & Coordinate spline & 0.3761 & 0.1713 & 149/192 & 34 & [0.000, 0.000]\\
$1/2$ & Full quadratic & 0.6129 & 0.3973 & 133/192 & 44 & [0.454, 0.928]\\
$1/2$ & Full cubic & 0.6660 & 0.4351 & 125/192 & 56 & [0.292, 0.612]\\
$1/2$ & Bilinear grid & 0.6158 & 0.3804 & 150/192 & 35 & [0.339, 0.674]\\
$3/4$ & Online & 0.4451 & 0.1496 & 192/192 & 0 & [0.000, 0.000]\\
$3/4$ & Quadratic (rank 6) & 0.3799 & 0.1764 & 134/192 & 44 & [0.407, 0.896]\\
$3/4$ & ReLU (rank 6) & 0.4199 & 0.2126 & 162/192 & 25 & [0.285, 0.532]\\
$3/4$ & Equivalent hinge & 0.4088 & 0.1955 & 162/192 & 25 & [0.285, 0.532]\\
$3/4$ & Coordinate spline & 0.3839 & 0.1724 & 150/192 & 32 & [0.091, 0.199]\\
$3/4$ & Full quadratic & 0.6136 & 0.4030 & 134/192 & 44 & [0.407, 0.896]\\
$3/4$ & Full cubic & 0.6608 & 0.4158 & 126/192 & 56 & [0.405, 0.801]\\
$3/4$ & Bilinear grid & 0.6109 & 0.3704 & 138/192 & 44 & [0.126, 0.403]\\
\bottomrule\end{tabular}
\par\smallskip\begin{minipage}{.98\linewidth}\footnotesize
Times are medians of all seven fresh-process 32-state repetitions, in seconds. Omitted support/trial pairs and update events refer to all 64 catalogue states. Extra mean cost is budgeted minus full-safe actual original-policy cost; outward numerical intervals, not confidence intervals. The declared primary budget is $2^{-16}$; the original loss target remains $1/256$ and width target $1/8192$. A positive extra cost is foregone improvement, not concealed as equal accuracy. All acquisition is additional and reported separately. Identical actions give exactly zero contrast. Update events can exceed the number of distinct changed states.
\end{minipage}\end{table}

\begin{table}[htbp]\caption{Original-policy quality and standalone acquisition at the primary budget}\label{tab:quality110}
\centering\scriptsize
\begin{tabular}{@{}llrrrr@{}}\toprule
$\alpha$ & Arm & Mean original cost & Mean regret upper & Max action bound & Acquire sec\\\midrule
$1/2$ & Online & [0.6083740, 0.6083753] & 0.00000107 & 0.001277 & 0.0001\\
$1/2$ & Quadratic (rank 6) & [0.6083751, 0.6083764] & 0.00000216 & 0.003797 & 0.1078\\
$1/2$ & ReLU (rank 6) & [0.6083751, 0.6083763] & 0.00000211 & 0.004343 & 0.0874\\
$1/2$ & Equivalent hinge & [0.6083751, 0.6083763] & 0.00000211 & 0.004343 & 0.0878\\
$1/2$ & Coordinate spline & [0.6083745, 0.6083758] & 0.00000160 & 0.004021 & 0.0733\\
$1/2$ & Full quadratic & [0.6083751, 0.6083764] & 0.00000216 & 0.003798 & 0.0190\\
$1/2$ & Full cubic & [0.6083750, 0.6083763] & 0.00000209 & 0.003829 & 0.0194\\
$1/2$ & Bilinear grid & [0.6083747, 0.6083759] & 0.00000173 & 0.003520 & 0.0185\\
$3/4$ & Online & [0.6180187, 0.6180199] & 0.00000106 & 0.001542 & 0.0001\\
$3/4$ & Quadratic (rank 6) & [0.6180198, 0.6180210] & 0.00000214 & 0.003874 & 0.1069\\
$3/4$ & ReLU (rank 6) & [0.6180196, 0.6180209] & 0.00000199 & 0.004116 & 0.0837\\
$3/4$ & Equivalent hinge & [0.6180196, 0.6180209] & 0.00000199 & 0.004116 & 0.0839\\
$3/4$ & Coordinate spline & [0.6180194, 0.6180206] & 0.00000177 & 0.003874 & 0.0757\\
$3/4$ & Full quadratic & [0.6180198, 0.6180210] & 0.00000214 & 0.003874 & 0.0183\\
$3/4$ & Full cubic & [0.6180198, 0.6180210] & 0.00000218 & 0.003874 & 0.0189\\
$3/4$ & Bilinear grid & [0.6180192, 0.6180205] & 0.00000159 & 0.003126 & 0.0178\\
\bottomrule\end{tabular}
\par\smallskip\begin{minipage}{.98\linewidth}\footnotesize
All 64 states, equal weights. Regret is sharpened with a separately charged original-law reference bracket; no exact optimum is provided for free. Action-distance entries are outward bounds from original action curvature two, not measured errors against an exact actor. Acquisition includes paid labels, all fitting trials, independent numerical fitting audit and serialization. Hinge pays the same shared ReLU acquisition; larger polynomial and grid payloads are explicit resource-frontier controls. Cumulative process RSS is not mislabeled isolated arm memory. No population welfare ordering or calibrated adoption is claimed.
\end{minipage}\end{table}


The primary table displays actual economic sacrifice as well as saved computation. Across all arms and risks, the primary budget omits 2,386 of 3,072 possible extra support/trial pairs and accepts 549 update events, versus 670 for the full safeguard. These are events, not necessarily distinct states. The largest upper endpoint of the mean additional-cost interval against the full safeguard is below $9.28\times10^{-7}$. All un-stopped final actions agree exactly with the full controls in the 64 arm--risk--budget cases; every stopped row satisfies its recorded ceiling. This verifies the comparison's trace premise rather than assuming that a smaller batch leaves floating arithmetic unchanged.

For ReLU at $\alpha=1/2$, the primary rule omits 153 of 192 pairs, reduces original terminal-risk queries from 94,737 to 35,832, and lowers the median 32-state clock from $0.4074$ to $0.2074$ seconds. Its original mean cost lies in $[0.6083751,0.6083763]$. The extra mean cost relative to the full safeguard is enclosed by $[1.95,3.34]\times10^{-7}$, while its mean reduction relative to its completed plain incumbent is enclosed by $[0.00017631,0.00017693]$. At $\alpha=3/4$, the corresponding query count falls from 95,251 to 32,881, and the median falls from $0.4199$ to $0.2126$ seconds. The extra mean cost is enclosed by $[2.85,5.32]\times10^{-7}$. These are comparisons of actual original-policy costs, not differences of regret upper bounds.

The online service illustrates when improvement work is unnecessary: it accepts no safeguard changes, and the budget rule omits all 192 extra pairs at either risk level while returning the same policy. Its medians fall from $0.4251$ to $0.1486$ and from $0.4451$ to $0.1496$ seconds in the new matched campaign. Several learned arms still cost more than that accurate online control after refinement. Quadratic and spline representations benefit from the same rule; their conventional victories are not removed. ReLU and equivalent hinges have exactly equal rows and trial records in all twelve risk-by-service comparisons. Their different clocks are implementation cost, not a separation of representable policies or acquired information.

The separate unchanged-source reproduction also supports the compiler as a common original-model improvement. Its online medians decrease from approximately $0.6395$ to $0.1380$ seconds at $1/2$ and from $0.6394$ to $0.1514$ at $3/4$. These clocks belong to that reproduction, not to the new six-service campaign. Both sets of seven raw timing repetitions are retained; one is not substituted for the other. The compiler does not change the deployed terminal controller, and its feasible witness remains the authority for continuation information.

\subsection{Error budgets, work accounting and adoption}
An action-distance bound is reported from original action curvature two and the verified regret bound. It is not measured error against a freely supplied exact actor. A separately charged 4,096-bin original-model reference bracket sharpens the reported regret account; its approximately $10.31$ seconds are benchmark work, not a free deployed-service oracle. A state exhausting all three attempts without stopping may have regret above $\theta$, but must still satisfy the original $1/256$ target. The stopping budget bounds foregone gain at stopped states; it is never relabeled the final regret target for every row.

Every acquisition charge includes the common original labels, all feature trials and exclusions, fitting, independent numerical audit, serialization and reload. Standalone comparisons charge a fitted service its full share of those costs; the complete research tournament additionally pays for all competing arms and budgets, reference work, all audits and all timing processes. The unchanged-source campaign took about $332.27$ seconds internally and the new campaign about $275.41$ seconds on the recorded host, before separately recorded verification and publication commands. Human research and code-development effort are not assigned fictitious zero economic cost.

All clocks are descriptive: seven sequential fresh processes, a fixed within-process order and one computational thread, with processor frequency and unrelated load not claimed controlled. We retain every repetition, operation category and cumulative process RSS. Query-event counts are not summed and called floating-point operations; cumulative process RSS is not an isolated per-arm memory allocation. The sensitivity appendix reports every declared budget, including larger cost sacrifices and smaller work savings. No theta or favorable repetition was selected after seeing outcomes.

The result is a certified computation--policy trade-off, not a calibrated purchase decision. A deployment would need an economically justified gain budget, workload and compute price, and an acquisition policy chosen without hindsight. The observed saving does not by itself establish neural adoption over a strong conventional method.

\subsection{Independent checks and retained adverse evidence}
The unchanged-source campaign retains 4,096 primary decisions and 14,336 timing decisions; its 18,432 rows and 27,648 safeguard-trial records pass 169,472 independent exact-arithmetic checks. The new campaign retains 6,144 primary decisions and 21,504 timing decisions; its 27,648 rows and all 69,120 active or omitted trial slots pass 794,058 checks. The latter includes 32 exact paired inherited-control identities and all 64 budget comparisons. A separate original-kernel audit passes 3,072 rational checks over 384 contexts at $d=2,4,8$. This is kernel scaling, not an unexecuted learned high-dimensional or longer-horizon tournament. Sixteen unchanged-source and twelve new preflight tests also pass. Finite tests support implementation evidence but do not replace the complete mathematical proofs.

The historical R97 campaign remains distinct: all 36 residual spans decreased, but total first-proposal rejections increased from 644 to 757. Quadratic rejections rose from 91 to 210, ReLU from 286 to 287, and spline fell from 267 to 260. Quadratic least squares was the observed complete-work winner. Its 1,800-cell sensitivity account had 156 adoption, 199 rejection and 1,445 unresolved cells, with zero learned-service adoption over quadratic least squares. All twelve retained R96 simultaneous policy-cost intervals cross zero. Neither the new stopping rule nor the restored tail theorem retroactively turns those experiments into successful action-aligned training.

A full inherited-test replay was attempted as a separate archive audit. It ran 435 cases across 21 suites, of which ten suites passed; twenty errors and one failed pinned-file assertion remain because old kernels or historical evidence paths are absent from the delivered layout. The actual logs are preserved. This is not an all-green full-history replay, and the successful current campaign cannot repair it by changing labels. The package separates the current executable dependency closure, verified outputs and mathematical claims from the outstanding historical archive dependencies.


\section{Conclusion}\label{sec:conclusion110}
Neural Bellman Operators concern continuation learning and feasible decisions under the original recursive economic objective. This revision connects their construction to action-aligned tail accuracy, original-model Bellman structure, completed-incumbent improvement and paid stopping. The remaining-gain rule has an explicit economic meaning: omitted computation can sacrifice at most a verified amount of original-policy improvement. Its recursive extension requires a comparator-future residual and risk-distorted propagation, not an unqualified average of physical state frequencies.

The original investment comparison executes that trade-off with common information and strong conventional controls. It records substantial savings in support and certification work, strictly improved fitted policies relative to completed plain incumbents, and the small but sometimes positive cost of stopping earlier than the full safeguard. Exact ReLU--hinge policy equivalence and the competitiveness of quadratic, spline and semi-analytic methods remain visible. No internal fitting metric or favorable timing subset is substituted for original economic loss.

The research program continues to target neural-specific representation and learning advantages under equal information and resources, prospective validation of tail-aware acquisition, genuine longer-horizon and multiple-control continuous-risk comparisons, and independently calibrated economic adoption. The recursive-utility BSDE bridge makes one rigorous connection to NDU/FBSDE research without transferring results across objectives or information structures. Current successful computations and complete proofs do not close those broader claims, nor do the preserved negative results replace their pursuit. The original NBO identity and the economic and numerical-methods ambition are retained.

\bibliographystyle{ecta-fullname}
\bibliography{references99,references110-additions}
\end{document}
